\subsection{Completely intersecting families}

Let A be a ring, M an A-module, \(x = (x_{i})_{i \in I}\) a family of elements of A.

DEFINITION 2. --- The family x is said to be completely secant for M if one has \(\mathrm{H}_{i}(\mathbf{x},\mathbf{M})=0\) for i\textgreater0.

If I is finite, it amounts to the same thing (X, p.~150) to say that one has \(\mathbf{H}^i (\pmb {x},\mathbf{M}) = 0\) for \(i <   \mathrm{Card}\) (I).

The following proposition allows us to give examples of completely secant families.

PROPOSITION 5. — Let \(x = (x_1, \ldots, x_n)\) be a sequence of elements of A. If for \(i=1, \ldots, n\) the homothety of ratio \(x_i\) in the A-module \(M/(x_1 M + \cdots + x_{i-1} M)\) is injective, the sequence \(x\) is completely secant for M.

A sequence satisfying the conditions of the proposition is said to be regular for M, or M-regular. Note that this property generally depends on the order of the \(x_i\) . For example, the sequence (1, 0) is always M-regular, whereas the sequence (0, 1) is M-regular only if M is zero. On the other hand, whether a sequence is completely secant does not depend on the order of the terms.

Let us prove the proposition by induction on \(n\) , the case \(n=0\) being immediate. Set \(x'=(x_{1},...,x_{n-1})\) ; if the sequence \(x\) is M-regular, the sequence \(x'\) is M-regular and multiplication by \(x_{n}\) to \(\mathrm{M}/(\mathbf{x}')\mathrm{M}\) is injective; by the induction hypothesis, we have \(\mathrm{H}_{i}(\mathbf{x}',\mathrm{M})=0\) for \(i>0\) ; it then follows from cor. 3 of X, p.~157, that \(\mathrm{H}_{i}(\mathbf{x},\mathrm{M})=0\) for \(i>0\) .

Example. --- Let \(k\) be a ring; take \(\mathbf{A} = k[\mathbf{X}_1, \ldots, \mathbf{X}_n]\) and \(\boldsymbol{x} = (\mathbf{X}_1, \ldots, \mathbf{X}_n)\) . The sequence \(\boldsymbol{x}\) is \(\mathbf{A}\) -regular and the proposition gives back acyclicity in degrees \(>0\) of the complex \(\mathbf{S}_k(k^n) \otimes_k \boldsymbol{\Lambda}_k(k^n)\) (cf.~X, p.~152, prop. 3).

Similarly in the ring of formal power series \(\hat{\mathbf{A}} = k[[\mathbf{X}_1, \ldots, \mathbf{X}_n]]\) , the sequence \((\mathbf{X}_1, \ldots, \mathbf{X}_n)\) is \(\hat{\mathbf{A}}\) -regular, hence completely secant for \(\hat{\mathbf{A}}\) .

PROPOSITION 6. --- a) If \(\sum_{i\in I}x_i\mathbf{A} = \mathbf{A}\) , the family \((x_i)_{i\in I}\) is completely secant for M.

\begin{enumerate}
\def\labelenumi{\alph{enumi})}
\setcounter{enumi}{1}
\tightlist
\item
  Let \(x = (x_1, \ldots, x_n)\) be a sequence of elements of A. Let \((a_{ij}) \in \mathbf{GL}_n(\mathbf{A})\) ; set
\end{enumerate}

\[
y _ {i} = \sum_ {j} a _ {i j} x _ {j}.
\]

If the sequence x is completely secant for M, so is the sequence \((y_{1}, \ldots, y_{n})\) . c) Let \(k_{1}, \ldots, k_{n}\) be integers \(\geqslant 1\) ; for the sequence \((x_{1}^{k_{1}}, \ldots, x_{n}^{k_{n}})\) to be completely secant for M, it is necessary and sufficient that the sequence x be completely secant for M. Assertion a) follows from cor. 1, p.~148.

Let us prove \(b\) ). Let \(f: \mathbf{A}^n \to \mathbf{A}^n\) be the isomorphism defined by the matrix \(^t(a_{ij})\) ; it follows from X, p.~151, that \(1_{\mathbf{M}} \otimes \Lambda(f)\) is an isomorphism of the complex \(\mathbf{K}.(y, \mathbf{M})\) on the complex \(\mathbf{K}.(x, \mathbf{M})\) , whence \(b\) .

To prove \(c\) ), it suffices evidently to show that if \(k\) is an integer \(\geqslant 1\) , the sequence \((x_{1}, \ldots, x_{n-1}, x_{n}^{k})\) is completely secant for M if and only if the same holds for the sequence \(x\) . Let \(x' = (x_{1}, \ldots, x_{n-1})\) ; by cor. 3, p.~157, the first condition (resp. the second) means that the homothety of ratio \(x_{n}^{k}\) (resp. \(x_{n}\) ) is bijective in H \(_{i}(x', M)\) for \(i \geqslant 1\) and injective in M/(x') M. These two conditions are clearly equivalent, whence \(c\) ).

Remark. --- 1) The assertion analogous to c) for regular sequences is true (X, p.~207, exercise 5).

PROPOSITION 7. --- a) Let N be a flat A-module. If the family x is completely secant for M, it is so for M ⊗\_A N.

\begin{enumerate}
\def\labelenumi{\alph{enumi})}
\setcounter{enumi}{1}
\tightlist
\item
  Let \(0 \to \mathbf{M}' \to \mathbf{M} \to \mathbf{M}'' \to 0\) be an exact sequence of A-modules. If the family \(x\) is completely secant for \(\mathbf{M}'\) and for \(\mathbf{M}''\) , it is so for \(\mathbf{M}\) .
\end{enumerate}

The complex \(\mathbf{K}.\left(\mathbf{x},\mathbf{M}\otimes_{\mathbf{A}}\mathbf{N}\right)\) is isomorphic by definition to \(\mathbf{K}.\left(\mathbf{x},\mathbf{M}\right)\otimes_{\mathbf{A}}\mathbf{N}\) ; since N is flat, we deduce an isomorphism H. \((\mathbf{x},\mathbf{M})\otimes_{\mathbf{A}}\mathbf{N}\to\mathrm{H}.\left(\mathbf{x},\mathbf{M}\otimes_{\mathbf{A}}\mathbf{N}\right)\) (X, p.~66, cor. 2), whence \(a\) ).

The complex \(\mathbf{K}.\left(\mathbf{x}\right)\) being flat, we have an exact sequence of complexes

\[
0 \rightarrow \mathbf {K}. (x, \mathrm{M} ^ {\prime}) \rightarrow \mathbf {K}. (x, \mathrm{M}) \rightarrow \mathbf {K}. (x, \mathrm{M} ^ {\prime \prime}) \rightarrow 0;
\]

assertion b) follows from the associated exact homology sequence.

Remarks. --- 2) The assertions analogous to a) and b) for regular sequences are immediate.

\begin{enumerate}
\def\labelenumi{\arabic{enumi})}
\setcounter{enumi}{2}
\tightlist
\item
  If the family \(x\) is completely secant for \(A\) , the complex \(K.(x, A)\) defines a free resolution of the \(A\) -module \(A/x\) , with \(x = \sum_{i \in I} x_i A\) ; we therefore have for every integer \(i \geqslant 0\) and for every \(A\) -module \(M\) isomorphisms
\end{enumerate}

\[
\operatorname{Ext} _ {\mathbf {A}} ^ {n - i} (\mathbf {A} / \mathfrak {x}, \mathbf {M}) \rightarrow \mathrm{H} ^ {n - i} (\mathfrak {x}, \mathbf {M}) \rightarrow \mathrm{H} _ {i} (\mathfrak {x}, \mathbf {M}) \rightarrow \operatorname{Tor} _ {i} ^ {\mathbf {A}} (\mathbf {A} / \mathfrak {x}, \mathbf {M}).\tag{15}
\]

\begin{enumerate}
\def\labelenumi{\arabic{enumi})}
\setcounter{enumi}{3}
\tightlist
\item
  We say that the sequence \(\boldsymbol{x} = (x_1, \dots, x_n)\) is M-coregular if (denoting by \((x_i)_M\) the homothety with ratio \(x_i\) in M) the homothety of ratio \(x_i\) in the module
\end{enumerate}

\[
\operatorname{Ker} \left(x _ {1}\right) _ {\mathbf {M}} \cap \dots \cap \operatorname{Ker} \left(x _ {i - 1}\right) _ {\mathbf {M}}
\]

is surjective for \(i = 1, \ldots, n\) . We then have \(\mathbf{H}_{i}(\mathbf{x}, \mathbf{M}) = 0\) for \(i < n\) ; this is proved in the same manner as prop. 5.

\[
\mathrm{A} = k \left[ \mathrm{D} _ {1}, \dots , \mathrm{D} _ {n} \right],
\]

\[
\mathbf {Q} \text {-algebra},
\]

\(M = k[T_1, \ldots, T_n]\) , endowed with the structure of an A-module such that \(D_i P = \partial P / \partial T_i\) for all \(P \in M (1 \leqslant i \leqslant n)\) . One verifies immediately that the sequence \((D_1, \ldots, D_n)\) is M-coregular; taking into account the example on p.~155, one deduces that the de Rham complex of \(k[T_1, \ldots, T_n]\) over \(k\) is acyclic in degrees \textgreater{} 0.

